\documentclass[10pt,a4paper]{amsart}
\usepackage[margin=19mm]{geometry}
\usepackage{amsmath,amssymb,mathtools}
\usepackage[scr=rsfs,cal=euler]{mathalfa}
\usepackage{xfrac}
\usepackage[unicode,hidelinks,bookmarks=false]{hyperref}
\newcommand{\ie}{i.e.\ }
\newcommand{\cf}{cf.\ }
\newcommand{\resp}{respectively\ }
\newcommand{\Subsection}{\subsection}
\providecommand{\nobreakdash}{\nobreak\hbox{-}}
\setlength{\emergencystretch}{2em}
\multlinegap0pt
\usepackage{booktabs}
\usepackage{amssymb}
\usepackage{mathtools}
\newtheorem{lemma}{Lemma}
\newtheorem{proposition}{Proposition}
\newcommand{\E}{\mathbb E}
\newcommand{\HS}{\mathrm{HS}}
\newcommand{\R}{\mathcal R}
\newcommand{\ip}[2]{\langle#1,#2\rangle}
\newcommand{\norm}[1]{\left\lVert#1\right\rVert}
\newcommand{\op}{\mathrm{op}}
\title{Resolvent characteristics and quadratic mixing\\of Kac's walk on $\mathrm{SO}(n)$}
\author{Yunjiang Jiang}
\date{Source: arXiv:2609.08018v1 (CC BY 4.0)}
\begin{document}
\maketitle

\noindent\emph{Source and licence.} Resolvent characteristics and quadratic mixing\\of Kac's walk on $\mathrm{SO}(n)$. Version arXiv:2609.08018v1 (CC BY 4.0), distributed by arXiv under the Creative Commons Attribution 4.0 International licence, http://creativecommons.org/licenses/by/4.0/, as recorded in the arXiv licence metadata for this exact version and shown on the abstract page https://arxiv.org/abs/2609.08018v1. Full licence notice: \texttt{licenses/arxiv-2609.08018-CC-BY-4.0.txt}.\par\smallskip

\noindent\textit{Editorial scope.} Counted unit: the article's proposition prop:partial-resolvent (source0.tex lines 805--815). The article's complete original proof of it is delivered with it (source0.tex lines 816--983, one \texttt{proof} environment of 168 lines). No mathematics is written, repaired or re-derived: the statement and the proof are the article's own lines, in the article's own notation, and no retained line is substituted or re-flowed.  Retained uncounted as the source-local context the proof needs: lines 352--364; lines 381--386; lines 498--501; lines 594--600; lines 679--690; lines 693--698; lines 708--712.  Nothing was omitted from any retained span, so the package carries no omission record.  No retained line contains a citation, so the wrapper carries no bibliography and no bibliography entry.  No retained line contains a figure environment, an image-inclusion command or drawing code, so no image is delivered and the package declares no asset dependency.  Editorial formatting only, in addition: an AMS \texttt{amsart} preamble that restates the article's own declarations and macros used by the retained lines, namely the environments \texttt{lemma}, \texttt{proposition} on separate counters, together with the standard packages those lines need.  The retained environments print in this document as Lemma 1, Proposition 1 in order of appearance, whereas the article prints the same items with the numbering of its own full document.  The complete native source of the article as distributed in the arXiv e-print for this exact version is bundled unchanged as \texttt{sources/209688/source0.tex}.
\begin{lemma}[Accumulating centered increments before a rotation]\label{lem:rotated-energy}
Let $\mathcal H$ be a finite-dimensional Hilbert space. Suppose
\[
 Z_{t+1}=\mathcal O_t(Z_t+b_t+\xi_t),
\]
where $\mathcal O_t$ is an orthogonal operator, $Z_t,b_t$ are measurable with respect to the past $\mathcal F_t$, and $\E(\xi_t\mid\mathcal F_t)=0$. The fresh orthogonal operator may depend on $\xi_t$. For $t_0<T$,
\begin{equation}\label{eq:rotation-energy-lemma}
 \|Z_T\|_{L^2(\mathcal H)}
 \le \|Z_{t_0}\|_{L^2(\mathcal H)}
 +\sum_{t=t_0}^{T-1}\|b_t\|_{L^2(\mathcal H)}
 +\left(\sum_{t=t_0}^{T-1}\E\|\xi_t\|_{\mathcal H}^2\right)^{1/2}.
\end{equation}
\end{lemma}

\begin{equation}\label{eq:Rbounds}
 \norm{\R_0 A}_F\le C\sqrt n\norm A_{\HS}
 \le Cn^{3/2}\norm A_{\op},\qquad
 \R(uv^*)=-2uv,\qquad
 \norm{\R_0(uv^*)}_F\le2\norm u\norm v.
\end{equation}

\begin{equation}\label{eq:parameters}
 \ell=\log(en),\qquad R=\lceil n^{5/4}\rceil,\qquad
 n^{-1/100}\le\eta\le1.
\end{equation}

\begin{lemma}\label{lem:diag-resolvent}
For a total colored covariance $S_t$, $R\le t\le4N$, and deterministic $z\in[\eta,5]$,
\begin{equation}\label{eq:diag-resolvent}
 \E V((S_t+zI)^{-1})\le C n^{-3/4}\eta^{-2}.
\end{equation}
With a norm weight $M^j$, $j\le8$, the right side may be multiplied by $C_j\ell^j$.
\end{lemma}

\begin{proposition}[Scalar stability]\label{prop:scalar}
Uniformly over these schedules,
\begin{equation}\label{eq:scalar-stability}
 \left\|\max_{t\le m}|s_t-\sigma|\right\|_{L^2}
 \le C n^{-3/8}\eta^{-4}.
\end{equation}
In particular, for an uncolored $2N$-step core,
\begin{equation}\label{eq:scalar-final}
 \E\tau(D_{2N}+\eta I)^{-1}\le2
\end{equation}
for sufficiently large $n$, uniformly in~\eqref{eq:parameters}.
\end{proposition}

\begin{equation}\label{eq:exact-resolvent}
 T_{t+1}^{\mathrm{pb}}
 =\widehat T-\frac{(\widehat T a)(\widehat T a)^*}
                       {1+\ip a{\widehat T a}},\qquad
 \widehat T=(S_t+(z-d)I)^{-1}.
\end{equation}

\begin{equation}\label{eq:rank-derivative-expanded}
 \partial_zQ_z
 =-\frac{(T_zt_a)t_a^*+t_a(T_zt_a)^*}{1+x}
   +\frac{y\,t_at_a^*}{(1+x)^2}.
\end{equation}

\begin{proposition}\label{prop:partial-resolvent}
Uniformly for every deterministic schedule of length at most $4N$ and every allowed terminal regularization $\eta$,
\begin{equation}\label{eq:partial-resolvent}
 \sup_{t\le m}\norm{\R_0T_t}_{L^2(F)}
 \le Cn^{5/4}\eta^{-6},\qquad
 \sup_{t\le m}\norm{\R_0(Y_tT_t)}_{L^2(F)}
 \le Cn^{5/4}\ell\eta^{-6}.
\end{equation}
In particular, for a colored covariance at any terminal time and fixed argument $\eta$, both squared partial traces have expectation at most
$Cn^{5/2}\ell^2\eta^{-12}$.
\end{proposition}

\begin{proof}
We display the error bounds in the two recurrences. At an active step let $T=T_t$, $Y=Y_t$, $U=YT$, $d=d_0$, and $\chi=\chi_t$. Pull back the common rotation. Besides~\eqref{eq:exact-resolvent}, the exact mixed update is
\begin{equation}\label{eq:mixed-exact}
 U_{t+1}^{\mathrm{pb}}
 =Y\widehat T+
   \frac{(\chi a-Y\widehat T a)(\widehat T a)^*}
        {1+\ip a{\widehat T a}}.
\end{equation}
Define
\[
 Q_T(a)=\frac{(Ta)(Ta)^*}{1+\ip a{Ta}},\qquad
 Q_U(a)=\frac{(\chi a-Ua)(Ta)^*}{1+\ip a{Ta}}.
\]

The mixed identity can be checked without any commutation assumption on $Y$ and $T$. If $\widehat x=\langle a,\widehat Ta\rangle$, then
\begin{align*}
 (Y+\chi aa^*)\left(\widehat T-
 \frac{(\widehat Ta)(\widehat Ta)^*}{1+\widehat x}\right)
 &=Y\widehat T+\chi a(\widehat Ta)^*
   -\frac{(Y\widehat Ta+\chi\widehat x a)(\widehat Ta)^*}{1+\widehat x}\\
 &=Y\widehat T+
   \frac{(\chi a-Y\widehat Ta)(\widehat Ta)^*}{1+\widehat x}.
\end{align*}
In particular, the ordered product in the leading term is $UT=YT^2$, not $TYT$.

Expanding the deterministic parameter shift gives
\begin{align}
 T_{t+1}^{\mathrm{pb}}&=T+dT^2-Q_T(a)+E_{T,a},\label{eq:Texp}\\
 U_{t+1}^{\mathrm{pb}}&=U+dUT+Q_U(a)+E_{U,a}.\label{eq:Uexp}
\end{align}
The remainders satisfy the following sufficient bounds:
\begin{align}
 \norm{\R_0E_{T,a}}_F&\le C N^{-1}\eta^{-4},\label{eq:ET}\\
 \norm{\R_0E_{U,a}}_F&\le C N^{-1}(1+\norm Y_{\op})\eta^{-4}.
 \label{eq:EU}
\end{align}
Here and below they are zero at inactive steps. To verify these bounds, the full-matrix Taylor remainder of $\widehat T$ has norm at most $d^2\eta^{-3}$; apply~\eqref{eq:Rbounds}, multiplying by $\norm Y$ in the mixed case. The rank-one terms and their derivatives with respect to the argument have dimension-free partial-trace bounds by~\eqref{eq:Rbounds}. For $Q_T$, the derivative is a sum of three rank-one terms with total product norm at most $C\eta^{-2}$, as in the preceding proof. For $Q_U$, differentiating its two numerator vectors and denominator bounds that total by $C(1+\norm Y)\eta^{-4}$. Integrate over an interval of length $d\le1/N$. Finally $n^{3/2}/N$ is bounded, so $d^2 n^{3/2}\eta^{-3}\le C N^{-1}\eta^{-4}$. This proves~\eqref{eq:ET}--\eqref{eq:EU}.

A detailed check for the mixed rank-one derivative is useful. At a variable argument $z$, write
\[
 Q_{U,z}(a)=\chi\frac{a(T_za)^*}{1+x}-YQ_z,
 \qquad x=\langle a,T_za\rangle.
\]
The derivative of the first term is a sum of two rank-one operators:
\[
 -\chi\frac{a(T_z^2a)^*}{1+x}
 +\chi\frac{y\,a(T_za)^*}{(1+x)^2},\qquad y=\|T_za\|^2.
\]
Its sum of product norms is at most $C\eta^{-3}$ for $\eta\le1$. Multiplication of the three rank-one derivatives in~\eqref{eq:rank-derivative-expanded} by $Y$ costs at most $\|Y\|$. Applying the dimension-free partial-trace bound separately to these five terms gives
\[
 \|\partial_z\R_0Q_{U,z}(a)\|_F
 \le C(1+\|Y\|)\eta^{-4}.
\]
For the deterministic Taylor remainders the general partial-trace inequality gives
$Cn^{3/2}d^2\eta^{-3}$, or this times $\|Y\|$. Integrating the rank-one derivatives over length $d\le N^{-1}$ proves the two claimed error bounds with an absolute constant. This use of rank-one norms, rather than the general partial-trace inequality on the random update, preserves the needed power of $n$.


Write $s=\tau T$ and
\[
 \delta_a=\frac1{1+\ip a{Ta}}-\frac1{1+s}.
\]
Then $N^{-1}\sum_a\delta_a^2\le V(T)$. The exact conditional means of the rank-one terms at fixed $T$ are
\begin{align}
 \E_aQ_T(a)&=\frac{T^2}{N(1+s)}+
                \frac1N T\operatorname{diag}(\delta_a)T,
                \label{eq:Tmean}\\
 \E_aQ_U(a)&=\frac{\chi T-UT}{N(1+s)}+
                \frac1N(\chi I-U)\operatorname{diag}(\delta_a)T.
                \label{eq:Umean}
\end{align}
By the general partial-trace bound, their last terms obey

For example, $\|\operatorname{diag}(\delta_a)\|_{\HS}=\sqrt{N\,N^{-1}\sum_a\delta_a^2}$. Hence
\[
 \frac{C\sqrt n}{N}\|T\|_{\op}^2
       \|\operatorname{diag}(\delta_a)\|_{\HS}
 \le C\sqrt{\frac nN}\eta^{-2}\sqrt{V(T)}.
\]
For the mixed term, use $\|\chi I-U\|\le1+\|Y\|/\eta\le(1+\|Y\|)/\eta$. These are the factors behind the following bounds:

\begin{align}
 \left\|\R_0\left[N^{-1}T\operatorname{diag}(\delta_a)T\right]\right\|_F
 &\le C\sqrt{n/N}\,\eta^{-2}\sqrt{V(T)},\label{eq:den-T}\\
 \left\|\R_0\left[N^{-1}(\chi I-U)\operatorname{diag}(\delta_a)T\right]\right\|_F
 &\le C\sqrt{n/N}(1+\norm Y)\eta^{-2}\sqrt{V(T)}.
 \label{eq:den-U}
\end{align}
For $t\ge R$, Lemma~\ref{lem:diag-resolvent} and its weighted form make their $L^2$ norms at most
\begin{equation}\label{eq:den-L2}
 Cn^{-7/8}\eta^{-3},\qquad C\ell n^{-7/8}\eta^{-3},
\end{equation}
respectively.

The leading matrix coefficient after combining the shift and the mean update is
\begin{equation}\label{eq:gcoef}
 g_t=d_0-\frac1{N(1+s_t)}
     =\frac{d_0(s_t-\sigma)}{1+s_t}.
\end{equation}
Therefore the unrotated mean increment of $\R_0T$ is
$g_t\R_0(T_t^2)$, apart from the denominator and Taylor errors above. Since
$\norm{\R_0(T_t^2)}_F\le Cn^{3/2}\eta^{-2}$,
Proposition~\ref{prop:scalar} gives
\begin{equation}\label{eq:leading-T}
 \norm{g_t\R_0(T_t^2)}_2
 \le \frac C N n^{9/8}\eta^{-6}.
\end{equation}
For the mixed resolvent, the corresponding increment is
\begin{equation}\label{eq:leading-U-form}
 g_t\R_0(U_tT_t)+\frac{\chi_t}{N(1+s_t)}\R_0T_t.
\end{equation}
The first term has $L^2$ norm at most
$C N^{-1}n^{9/8}\ell\eta^{-6}$.
Indeed $\norm{\R_0(U_tT_t)}\le Cn^{3/2}\norm{Y_t}\eta^{-2}$, and the cutoff rule bounds
$\norm{\norm{Y_t}|s_t-\sigma|}_2$ by
$C\ell n^{-3/8}\eta^{-4}$ plus a negligible error.
The second term is at most $N^{-1}\norm{\R_0T_t}_2$ in $L^2$.

We explain how to sum these increments. For either partial trace, its exact recursion has the form
\[
 Z_{t+1}=V_t(Z_t+\beta_t+\xi_t)V_t^T,
 \qquad \E(\xi_t\mid\mathcal F_t)=0,
\]
where $V_t$ is the fresh fundamental rotation and $\beta_t$ is the conditional mean of the pulled-back increment. The rotation may depend on the centered increment. Define
$M_{t+1}=V_t(M_t+\xi_t)V_t^T$, starting from zero at time $R$. Norm invariance removes $V_t$ before taking expectations, so
\[
 \E\norm{M_T}_F^2=\sum_{t=R}^{T-1}\E\norm{\xi_t}_F^2.
\]
The difference $Z_T-M_T$ has norm at most
$\norm{Z_R}+\sum_{t\ge R}\norm{\beta_t}$.
This argument needs no assertion that the rotated increments are martingale differences.

The exact random parts in~\eqref{eq:exact-resolvent} and~\eqref{eq:mixed-exact}, together with the rank-one bound, show that the centered-increment $L^2$ norms are at most
$C\eta^{-2}$ and $C\ell\eta^{-2}$ respectively. Their accumulated noise norms over $4N$ steps are therefore at most $Cn\eta^{-2}$ and $Cn\ell\eta^{-2}$.
For the first $R$ steps, simply sum the exact increments. The deterministic shift contributes at most $C d_0n^{3/2}\eta^{-2}$ per step, or this times $\norm{Y_t}$ in the mixed case, and the rank-one term contributes at most $C\eta^{-2}$ or $C(1+\norm{Y_t})\eta^{-2}$. Thus their initial $L^2$ norms are bounded by
\begin{equation}\label{eq:initial-matrix}
 CR\eta^{-2},\qquad CR\ell\eta^{-2}.
\end{equation}
This also covers any terminal time less than $R$.


For reference, the complete $L^2$ ledger for the unweighted partial trace is
\begin{center}
\small
\begin{tabular}{p{0.38\textwidth}p{0.23\textwidth}p{0.28\textwidth}}
\toprule
Contribution & Per active step, after time $R$ & Accumulated through $4N$\\
\midrule
Taylor remainder & $CN^{-1}\eta^{-4}$ & $C\eta^{-4}$\\
Denominator error & $Cn^{-7/8}\eta^{-3}$ & $Cn^{9/8}\eta^{-3}$\\
Scalar-error drift & $CN^{-1}n^{9/8}\eta^{-6}$ & $Cn^{9/8}\eta^{-6}$\\
Centered noise & $C\eta^{-2}$ & $Cn\eta^{-2}$\\
Initial interval & --- & $Cn^{5/4}\eta^{-2}$\\
\bottomrule
\end{tabular}
\end{center}
Only the noise row is summed quadratically; the three drift rows use Minkowski. Every row for the mixed partial trace has at most an additional factor $C\ell$, except its source $\chi\R_0T/[N(1+s)]$, which sums to at most $4\sup_t\|\R_0T_t\|_2$. Lemma~\ref{lem:rotated-energy} is the precise accumulation rule.

For $\R_0T$, summing~\eqref{eq:ET}, \eqref{eq:den-L2}, and~\eqref{eq:leading-T}, with the noise and initial bounds, gives
\[
 \norm{\R_0T_T}_2\le C\left(
 n^{5/4}\eta^{-2}+n\eta^{-2}
 +n^{9/8}\eta^{-6}+n^{9/8}\eta^{-3}+\eta^{-4}\right)
 \le Cn^{5/4}\eta^{-6}.
\]
The bound is uniform in $T\le m$. Use it in the second term of~\eqref{eq:leading-U-form}. Its accumulation is at most four times that uniform bound. All the other mixed estimates are the preceding ones with at most a factor $C\ell$. Consequently
$\norm{\R_0U_T}_2\le Cn^{5/4}\ell\eta^{-6}$.
This proves~\eqref{eq:partial-resolvent}. All discarded tails are smaller than the displayed bounds uniformly for $\eta\ge n^{-1/100}$.
\end{proof}

\end{document}
